% -*- mode: latex; TeX-master: t; -*-

% \documentclass[notes=hide]{beamer}
% \documentclass[notes=only]{beamer}
\documentclass[
14pt,
aspectratio=169,
usenames,
dvipsnames,
% handout,
x11names]{beamer}
%\let\Tiny=\tiny

% Font selection
\overfullrule=5pt

\usepackage{etex}
\usepackage{pgfpages}

\usepackage{tabularx}
\usepackage{multicol}
% \usepackage[subfolder]{gnuplottex}

%% beamerscape dependencies
\usepackage[absolute,overlay]{textpos}
\setlength{\TPHorizModule}{\paperwidth}
\setlength{\TPVertModule}{\paperheight}
\textblockorigin{0mm}{0mm}

\definecolor{greenvm}{HTML}{AFD18E}
\definecolor{g}{HTML}{0E3B12}
\definecolor{v}{HTML}{003366}

\newcommand{\src}[1]{\scriptsize Figure Source: \textit{#1}}

%%%%%%%%%%%%%%%%%%%%%%%%%%%%%%%%%%%%%%%%%%%%%%%%%%%%%%%%%%%%%%%%%%%%%%%%%%%%%%%%
% Beamer template
\usetheme[
	bullet=square,             % Default option: circle
%	nobigpagenumber,         % No big circled page number on lower right
	noheadline,		 % No headline
%       nofootline,		 % No footline
	% grey,                    % For a neutral grey background color
%	watermark=BG_both,	 % png file for the watermark
]{Ithaca}
% \usetheme{metropolis}
% \usetheme{Boadilla}
%\usetheme{Berlin}
% \usetheme{Warsaw}
% % \usetheme{Dresden}
% %\usecolortheme{seagull}
% \usecolortheme{beaver}
% %\usecolortheme{seahorse}
%\usefonttheme{serif}
% \usefonttheme{structurebold}

%%%%%%%%%%%%%%%%%%%%%%%%%%%%%%%%%%%%%%%%%%%%%%%%%%%%%%%%%%%%%%%%%%%%%%%%%%%%%%%%
% Template
%%%%%%%%%%%%%%%%%%%%%%%%%%%%%%%%%%%%%%%%%%%%%%%%%%%%%%%%%%%%%%%%%%%%%%%%%%%%%%%%

\definecolor{mypurple}{HTML}{494949}
\definecolor{myyellow}{HTML}{FFCC33}
\definecolor{mygrey}{HTML}{999999}

\setbeamercolor{alerted text}{fg=myyellow}
\setbeamercolor{headline}{bg=footline.bg}
\setbeamercolor{block body alerted}{bg=normal text.bg!90!black}
\setbeamercolor{block body}{bg=normal text.bg!90!black}
\setbeamercolor{block body example}{bg=normal text.bg!90!black}
\setbeamercolor{block title alerted}{fg=myyellow}
\setbeamercolor{block title}{bg=blue}
\setbeamercolor{block title example}{use={normal text,example text},fg=example text.fg!75!normal text.fg,bg=normal text.bg!75!black}
\setbeamercolor{frametitle}{fg=white}
\setbeamercolor{item projected}{fg=white}
\setbeamercolor{normal text}{bg=mypurple,fg=white}
\setbeamercolor{section in sidebar}{fg=brown}
\setbeamercolor{section in sidebar shaded}{fg=grey}
\setbeamercolor{separation line}{}
\setbeamercolor{sidebar}{bg=red}
\setbeamercolor{sidebar}{parent=palette primary}
\setbeamercolor{structure}{bg=black, fg=white}
\setbeamercolor{subsection in sidebar}{fg=brown}
\setbeamercolor{subsection in sidebar shaded}{fg=grey}

% \setbeamercolor{palette primary}{fg=mypurple!80!black}
% \setbeamercolor{palette secondary}{fg=mypurple!80!black}
% \setbeamercolor{palette quaternary}{fg=mypurple!50!black}
% \setbeamercolor{page number in foot}{fg=white}

% lets add page numbers
% \expandafter\def\expandafter\insertshorttitle\expandafter{%
%   \insertshorttitle\hfill\insertframenumber\,/\,\inserttotalframenumber}

% \setbeamercovered{transparent}
\beamertemplatetransparentcoveredhigh

% *****************************************
% >>>> 15.3.2    Using Beamer's Templates
% *****************************************
% As a user of the beamer class you typically do not 'use' or 'invoke'
% templates yourself, directly. For example, the frame title template
% is automatically invoked by beamer somewhere deep inside the frame
% typesetting process. The same is true of most other
% templates. However, if, for whatever reason, you wish to invoke a
% template yourself, you can use the following command.
% \usebeamertemplate***{ element name }
% -------------------------------
%%% 7.2.1 The Headline and Footline
% \setbeamertemplate{headline} % Beamer-Template/-Color/-Font
% \setbeamertemplate{headline}
% {%
%   \begin{beamercolorbox}{section in head/foot}
%     \vskip2pt\insertnavigation{\paperwidth}\vskip2pt
%   \end{beamercolorbox}%
% }
% \setbeamertemplate{headline}[default] % The default is just an empty headline.
% \setbeamertemplate{headline}[infolines theme]
% \setbeamertemplate{headline}[miniframes theme]
% \setbeamertemplate{headline}[sidebar theme]
% \setbeamertemplate{headline}[smoothtree theme]
% \setbeamertemplate{headline}[smoothbars theme]
% \setbeamertemplate{headline}[tree]
% \setbeamertemplate{headline}[split theme]
% \setbeamertemplate{headline}[text line]{ text } % The headline is typeset with 'text'
% Logo in every slide except title
% \addtobeamertemplate{headline}{}
% {%
% \vspace{1ex}
% \hspace{1ex}
% \includegraphics[height=1cm]{myrmigki_color}
% }
% -------------------------------
% \setbeamertemplate{footline} % Beamer-Template/-Color/-Font
% \setbeamertemplate{footline}[default]
% \setbeamertemplate{footline}[infolines theme]
% \setbeamertemplate{footline}[miniframes theme]
% \setbeamertemplate{footline}[page number]
% \setbeamertemplate{footline}[frame number]
% \setbeamertemplate{footline}[split]
% \setbeamertemplate{footline}[text line]{ text }
% % Footline in every slideexcept title
\setbeamertemplate{footline}{
  \leavevmode%
  \begin{beamercolorbox}[left,wd=.4\paperwidth,ht=2.5ex,dp=1.125ex,leftskip=.3cm, rightskip=.3cm plus1fil]{footline}%
    \insertshorttitle%
  \end{beamercolorbox}
  %
  \begin{beamercolorbox}[center,wd=.2\paperwidth,ht=2.5ex,dp=1.125ex,leftskip=.3cm,rightskip=.3cm plus1fil]{footline}%
    \centering
    \insertframenumber~of~\inserttotalframenumber%
  \end{beamercolorbox}%
  %
  \begin{beamercolorbox}[right,wd=.4\paperwidth,ht=2.5ex,dp=1.125ex,leftskip=.3cm,rightskip=.3cm plus1fil]{footline}%
    \hskip2ex plus1fill%
    \includegraphics[height=1.5ex]{cc}~%
    \includegraphics[height=1.5ex]{by}~~%
    % \includegraphics[height=1.5ex]{sa}~~%
    \insertshortauthor~%
  \end{beamercolorbox}%
}%
% -------------------------------
%%% 7.2.2 The Sidebars
% -------------------------------
%%% 7.2.3 Navigation Bars (funktioniert nur mit miniframe Themes)
% \setbeamertemplate{mini frames}[default] % shows small circles as mini frames.
% \setbeamertemplate{mini frames}[box] % shows small rectangles as mini frames.
% \setbeamertemplate{mini frames}[tick] % shows small vertical bars as mini frames.
% -------------------------------
%%% 7.2.4 The Navigation Symbols
%%% Beamer-Template/-Color/-Font navigation symbols
% \setbeamertemplate{navigation symbols}{} % suppresses all navigation symbols:
% \setbeamertemplate{navigation symbols}[horizontal] % Organizes the navigation symbols horizontally.
% \setbeamertemplate{navigation symbols}[vertical] % Organizes the navigation symbols vertically.
% \setbeamertemplate{navigation symbols}[only frame symbol] % Shows only the navigational symbol for navigating frames.
% -------------------------------
%%% 7.2.5 The Logo
% \setbeamertemplate{logo} % Beamer-Template/-Color/-Font
% -------------------------------
%%% 7.2.6 The Frame Title
% \setbeamertemplate{frametitle} % Beamer-Template/-Color/-Font
% \setbeamertemplate{frametitle}[default][center] % left, center, right
% \setbeamertemplate{frametitle}[shadow theme]
% \setbeamertemplate{frametitle}[sidebar theme]
% \setbeamertemplate{frametitle}[smoothbars theme]
% \setbeamertemplate{frametitle}[smoothtree theme]
% \setbeamertemplate{frametitle}
% {
%     \nointerlineskip
%     \begin{beamercolorbox}[sep=0.3cm,ht=2.5em,wd=\paperwidth]{frametitle}
%         \vbox{}\vskip-2ex%
%         \hspace{1.2cm}
%         \strut\insertframetitle\strut
%         \vskip-0.8ex%
%     \end{beamercolorbox}
% }
% % Minimalistic frametitle
% \setbeamertemplate{frametitle}
% {
%   \vspace{.2cm}
%   \hspace{-.8cm}
%   \insertframetitle
% }
% -------------------------------
%%% 7.2.7 The Background
% \setbeamertemplate{background canvas} % Beamer-Template/-Color/-Font
% \setbeamertemplate{background canvas}[default]
% \setbeamertemplate{background canvas}[vertical shading][ color options ] installs a vertically shaded background.
%     - top= color specifies the color at the top of the page. By default, 25% of the foreground of
%       the beamer-color palette primary is used.
%     - bottom= color specifies the color at the bottom of the page. By default, the background of
%       normal text at the moment of invocation of this command is used.
%     - middle= color specifies the color for the middle of the page. Thus, if this option is given, the
%       shading changes from the bottom color to this color and then to the top color.
%     - midpoint= factor specifies at which point of the page the middle color is used. A factor of 0
%       is the bottom of the page, a factor of 1 is the top. The default, which is 0.5 is in the middle.
% \setbeamertemplate{background canvas}[vertical shading][top=gray!40,bottom=gray!40,middle=white,midpoint=.9]
%\setbeamertemplate{background canvas}[vertical shading][top=white,bottom=gray!40]
% \setbeamercolor{background canvas}{bg=white}
%% Slide Background
\setbeamertemplate{background canvas}{%
	\begin{tikzpicture}
		\path [outer color = mypurple,
		       inner color = mypurple!75!white]
			(0,0) rectangle (\paperwidth,\paperheight);
	\end{tikzpicture}%
}

% \setbeamertemplate{background} % Beamer-Template/-Color/-Font
% \setbeamertemplate{background}[default] % is empty.
% \setbeamertemplate{background}[grid][step=1cm] % places a grid on the background.
%     - step= dimension specifies the distance between grid lines. The default is 0.5cm.
%     - color= color specifies the color of the grid lines. The default is 10% foreground.
% % Add global transition effect \transfade (breaks the background)
% \addtobeamertemplate{background canvas}{%
%   \transfade[duration=.1]%
% }{}
% -------------------------------
%%% 7.3 Margin Sizes
\setbeamersize{text margin left=2em,text margin right=2em}
% \setbeamersize{sidebar width left=2cm}
%         - text margin left= TEX dimension sets a new left margin. This excludes the left sidebar. Thus,
%           it is the distance between the right edge of the left sidebar and the left edge of the text.
%         - text margin right= TEX dimension sets a new right margin.
%         - sidebar width left= TEX dimension sets the size of the left sidebar. Currently, this command
%           should be given before a shading is installed for the sidebar canvas.
%         - sidebar width right= TEX dimension sets the size of the right sidebar.
%         - description width= TEX dimension sets the default width of description labels, see Section 11.1.
%         - description width of= text sets the default width of description labels to the width of the
%             text , see Section 11.1.
%         - mini frame size= TEX dimension sets the size of mini frames in a navigation bar. When two
%           mini frame icons are shown alongside each other, their left end points are TEX dimension far
%           apart.
%         - mini frame offset= TEX dimension set an additional vertical offset that is added to the mini
%           frame size when arranging mini frames vertically.
% -------------------------------
%%% 9.1 Adding a Title Page
% \setbeamersize{title page} % Beamer-Template/-Color/-Font
%    This template is invoked when the \titlepage command is used.
%    The following commands are useful for this template:
%     -  \insertauthor inserts a version of the author's name that is useful for the title page.
%     -  \insertdate inserts the date.
%     -  \insertinstitute inserts the institute.
%     -  \inserttitle inserts a version of the document title that is useful for the title page.
%     -  \insertsubtitle inserts a version of the document title that is useful for the title page.
%     -  \inserttitlegraphic inserts the title graphic into a template.
\defbeamertemplate*{title page}{flat2}[1][]
{
  \addtocounter{framenumber}{-1}
  \begin{center}
    \Huge{\inserttitle}\\[1ex]
    \hrule
    \vspace{1ex}
    \Large{\insertsubtitle}
    \vfill

    % \large\textcolor{normal text.fg}{\textit{\insertinstitute}, \insertdate}
    \vspace{1em}

    \vfill

    \large \insertauthor\\
    \vspace{-1ex}
    {\scriptsize \insertdate}\\
    \insertlogo
  \end{center}
}

% -------------------------------
%%% 9.2 Adding Sections and Subsections
% -------------------------------
%%% Parent Beamer-Template sections/subsections in toc
% This is a parent template, whose children are section in toc and subsection in toc.
% \setbeamertemplate{sections/subsections in toc}[default]
% \setbeamertemplate{sections/subsections in toc}[sections numbered]
% \setbeamertemplate{sections/subsections in toc}[subsections numbered]
% \setbeamertemplate{sections/subsections in toc}[circle]
% \setbeamertemplate{sections/subsections in toc}[square]
% \setbeamertemplate{sections/subsections in toc}[ball]
% \setbeamertemplate{sections/subsections in toc}[ball unnumbered]
% -------------------------------
%%% 9.6 Adding a Bibliography
% -------------------------------
% \setbeamertemplate{bibliography item} % Beamer-Template/-Color/-Font
% \setbeamertemplate{bibliography item}[default] %  little article icon as the reference
% \setbeamertemplate{bibliography item}[article] % Alias for the default.
% \setbeamertemplate{bibliography item}[book] % little book icon as the reference
% \setbeamertemplate{bibliography item}[triangle] % triangle as the reference
% \setbeamertemplate{bibliography item}[text] % reference text (like '[Dijkstra, 1982]')
% -------------------------------
%%% 10.1 Adding Hyperlinks and Buttons
% -------------------------------
%%% 11.1 Itemizations, Enumerations, and Descriptions
% \setbeamertemplate{items} % parent template of itemize items and enumerate items
% \setbeamertemplate{itemize items} % Parent Beamer-Template
% \setbeamertemplate{itemize items}[triangle]
% \setbeamertemplate{itemize items}[circle]
% \setbeamertemplate{itemize items}[square]
% \setbeamertemplate{itemize items}[ball]
% -------------------------------
% \setbeamertemplate{enumerate items}[default] % Numbered
% \setbeamertemplate{enumerate items}[circle] % Places the numbers inside little circles.
% \setbeamertemplate{enumerate items}[square] % Places the numbers on little squares.
% \setbeamertemplate{enumerate items}[ball] % 'Projects' the numbers onto little balls.
% -------------------------------
%%% 11.2 Hilighting
% -------------------------------
%%% 11.3 Block Environments
% \setbeamertemplate{blocks} % Parent Beamer-Template
% \setbeamertemplate{blocks}[default]
% \setbeamertemplate{blocks}[rounded][shadow=true]
 % \setbeamertemplate{blocks}[rounded][shadow=false]
% -------------------------------
%%% 11.4 Theorem Environments
% \setbeamertemplate{qed symbol} % Beamer-Template/-Color/-Font
% -------------------------------
% \setbeamertemplate{theorems} % Parent Beamer-Template
% \setbeamertemplate{theorems}[default]
% \setbeamertemplate{theorems}[normal font]
% \setbeamertemplate{theorems}[numbered]
% \setbeamertemplate{theorems}[ams style]
% -------------------------------
%%% 11.6 Figures and Tables
% \setbeamertemplate{caption} % Beamer-Template/-Color/-Font
% \setbeamertemplate{caption}[default] typesets the caption name (a word like 'Figure' or 'Abbildung' or 'Table')
% \setbeamertemplate{caption}[numbered] adds the figure or table number to the caption.
% \setbeamertemplate{caption}
% -------------------------------
% \setbeamertemplate{caption name} % Beamer-Color/-Font
% -------------------------------
%%% 11.10    Abstract
% -------------------------------
%%% 11.11 Verse, Quotations, Quotes
% -------------------------------
%%% 11.12 Footnotes
% -------------------------------
%%% 18.1 Specifying Note Contents
% \setbeamertemplate{note page} % Beamer-Template/-Color/-Font
% \setbeamertemplate{note page}[default]
% \setbeamertemplate{note page}[compress]
% \setbeamertemplate{note page}[plain]
% -------------------------------
%%% Specifying Which Notes and Frames Are Shown
 % \setbeameroption{hide notes}
% \setbeameroption{show notes}
% \setbeameroption{show notes on second screen= right }
% \setbeameroption{show only notes}


\renewcommand{\footnoterule}{}

\usepackage{hyperref}
\hypersetup{
  linktoc=all,
  bookmarks=false,           % show bookmarks bar?
  bookmarksopen,
  bookmarksnumbered,
  colorlinks = false,
  linkcolor=black,           % color of interlinks
  citecolor=black,           % color of the sitation links
  urlcolor=black,            % the url color
  unicode=true,              % non-Latin characters in Acrobat’s bookmarks
  pdftoolbar=true,           % show Acrobat’s toolbar?
  pdfmenubar=true,           % show Acrobat’s menu?
  pdffitwindow=true,         % window fit to page when opened
  pdftitle={Managed Runtime Systems},  % title
  pdfborder={ 0 0 0 },        % uBorder tin links
  pdfauthor = {Foivos Zakkak},
  pdfcreator = {Foivos Zakkak},
}

\usefonttheme{professionalfonts}% use own font handling
\usepackage{fontspec}
\usepackage{xunicode}
% \setmainfont{Comfortaa}
% \setsansfont{Comfortaa}
% \setsansfont{Liberation Sans}
\setmainfont[
% SmallCapsFont={Linux Biolinum},
SmallCapsFeatures={Letters=SmallCaps},
]{Liberation Sans}
% \setmonofont{Liberation Mono}
% \setmonofont{Inconsolata LGC for Powerline}
% Math fonts
\usepackage{unicode-math}
\setmathfont{xits-math.otf}

\usepackage{graphicx}
\usepackage{ulem}
\usepackage{color}
\usepackage{xspace}
\usepackage{xcolor}
\usepackage{array}
\usepackage{tikz}
\usetikzlibrary{arrows,arrows.meta,shapes,decorations,decorations.pathmorphing,calc,shadows,shadows.blur,shapes.multipart,positioning,patterns,fit,backgrounds,trees,shapes}
\usepackage{pgfplots}

\tikzset{
    %Define standard arrow tip
    >=stealth',
    % Define arrow style
    arrow/.style={
           ->,
           thick,
           color=uompurple,
           shorten <=2pt,
           shorten >=4pt,},
    % Daniel's proposal for "uncovering" parts of a tikz-tree %
    invisible/.style={opacity=0},
    visible on/.style={alt=#1{}{invisible}},
    alt/.code args={<#1>#2#3}{%
      \alt<#1>{\pgfkeysalso{#2}}{\pgfkeysalso{#3}} % \pgfkeysalso doesn't change the path
    },
  }
\pgfdeclaredecoration{penciline}{initial}{
    \state{initial}[width=+\pgfdecoratedinputsegmentremainingdistance,
    auto corner on length=1mm,]{
        \pgfpathcurveto%
        {% From
            \pgfqpoint{\pgfdecoratedinputsegmentremainingdistance}
                      {\pgfdecorationsegmentamplitude}
        }
        {%  Control 1
        \pgfmathrand
        \pgfpointadd{\pgfqpoint{\pgfdecoratedinputsegmentremainingdistance}{0pt}}
                    {\pgfqpoint{-\pgfdecorationsegmentaspect
                     \pgfdecoratedinputsegmentremainingdistance}%
                               {\pgfmathresult\pgfdecorationsegmentamplitude}
                    }
        }
        {%TO
        \pgfpointadd{\pgfpointdecoratedinputsegmentlast}{\pgfpoint{1pt}{1pt}}
        }
    }
    \state{final}{}
  }
\usepackage{amsmath}
\usepackage{textpos}
% \usepackage{subfigure}

\graphicspath{{figures/}}
% \newcommand{\longname}{\emph{Source-Level Compiler Optimizations for Task-Parallelism}}
\newcommand{\java}[0]{Java\texttrademark\xspace}
\newcommand{\code}[1]{\texttt{#1}}

\usepackage{bbding} % for the Checkmark and XSolidBrush
\newcommand{\tik}[0]{{\color{YellowGreen}\Checkmark}} % Check-mark
\newcommand{\ex}[0]{{\color{BrickRed}\XSolidBrush}}  % X-mark

% \usepackage{verbatim}
\usepackage{listings}
\lstset{
  language=c,
  basicstyle=\small\ttfamily,
%  columns=flexible,
  % numbers=left,
  % numberstyle=\ttfamily\tiny,
  showstringspaces=false,
  alsoletter={-},
  literate={-}{-}1,
%  numbersep=1em,
  xleftmargin=2em,
%  xrightmargin=1em,
  escapeinside={@}{@},
%  morecomment=[l][\color{BrickRed}]{\#pragma\ task},
  commentstyle=\color{black},
  morekeywords={final, transient, synthetic},
%   identifierstyle=\color{green},
%  backgroundcolor=\color{Honeydew1},
  keywordstyle=\bfseries\color{myyellow},
  stringstyle=\color{YellowGreen},
%  moredelim=**[is][\only<+(1)->{\color{black}\lstset{style=highlight}}]{~}{~},
}

\usepackage{booktabs} % bottomrule
\usepackage{multirow}

% Print table of contects at the beggining of each section
% \AtBeginSection[]{\begin{frame}\frametitle{Table of Contents}\tableofcontents[currentsection,currentsubsection]\end{frame}}
% \AtBeginSubsection[]{\begin{frame}\frametitle{Table of Contents}\tableofcontents[currentsection,currentsubsection]\end{frame}}

% Add outlines
% \AtBeginSection[]
% {
%   {
%   \setbeamertemplate{footline}{}
%   \begin{frame}<beamer>[plain,noframenumbering]
%     \frametitle{Outline}
%     \tableofcontents[currentsection]
%   \end{frame}
%   }
% }
\setcounter{tocdepth}{1}

%%%%%%%%%%%%%%%
% Intro slide %
%%%%%%%%%%%%%%%

\title{Managed Runtime Systems}
\subtitle{Lecture 05: Just In Time Compilation Part 1}
\author[\url{https://foivos.zakkak.net}]{Foivos Zakkak}
% \institute[]{ManLang'17}
\date{\url{https://foivos.zakkak.net}}
% \logo{\includegraphics[height=1.0cm]{uniman-logo}}
%% Spread those bullets
% \let\olditem\item
% \renewcommand{\item}{\setlength{\itemsep}{\fill}\olditem}

\begin{document}
\setbeamercovered{invisible}

% \maketitle

\begin{frame}[plain]
  \titlepage
  \centering
  \includegraphics[height=.75cm]{cc}~
  \includegraphics[height=.75cm]{by}\\[1em]
  % \includegraphics[height=.75cm]{nc-eu}~
  % \includegraphics[height=.75cm]{sa}\\[1em]
  \scriptsize{Except where otherwise noted, this presentation is licensed under the\\
    \href{http://creativecommons.org/licenses/by/4.0/}%
    {Creative Commons Attribution 4.0 International License.}\\[1ex]
    Third party marks and brands are the property of their respective
    holders.}
\end{frame}

% \begin{frame}[plain,noframenumbering]{Outline}
%     \tableofcontents
% \end{frame}

\begin{frame}{Acknowledgments}
  The following slides are based on the corresponding slides of Mario Walczko about \href{https://www.dropbox.com/s/hqc2axihgvxtadf/10\%20Dynamic\%20compilation.pdf}{\emph{Dynamic Compilation}}
\end{frame}

\section{Introduction}

\begin{frame}{Interpretation is Slow}
  \begin{itemize}  \setlength{\itemsep}{\fill}
  \item Step a single (or a few) bytecode(s) a time
  \item Not suitable execution patterns for HW-prefetchers
  \item Not enough data at hand to perform optimizations
  \end{itemize}
\end{frame}

\begin{frame}{Dynamic Compilation to the Rescue}
  \begin{itemize}  \setlength{\itemsep}{\fill}
  \item Produce machine code
  \item Better execution patterns for HW-prefetchers
  \item More data at hand to perform optimizations
  \item Generation of HW-specific code
  \item Generation of case-specific code
  \end{itemize}
\end{frame}


\section{Compiler Design Decisions}

\begin{frame}
  \centering
  {Dynamic Compiler Design Decisions}
\end{frame}

\begin{frame}{How Much to Compile}
  \begin{itemize}  \setlength{\itemsep}{\fill}
  \item Basic Block
  \item Method
  \item Multiple Methods (via inlining)
  \item Multiple Consecutive Methods (via tracing)
  \end{itemize}
\end{frame}

\begin{frame}{When to Compile}
  \begin{itemize}  \setlength{\itemsep}{\fill}
  \item Ahead of Time (AoT)
    \begin{itemize}
    \item At class loading
    \item At installation (see Android's dex2oat)
    \end{itemize}
  \item Just in Time (JIT)
    \begin{itemize}
    \item Compile at first invocation (stall till compilation completes)
    \item Re-use on later invocations
    \end{itemize}
  \item Practical JIT
    \begin{itemize}
    \item Start compiling after $n$ invocations
    \item Interpret while compiling to avoid stalls
    \end{itemize}
  \end{itemize}
\end{frame}

\begin{frame}{How Much to Optimize}
  \begin{itemize}  \setlength{\itemsep}{\fill}
  \item Minimal
    \begin{itemize}
    \item Replace bytecodes with function calls, like macro expansion
    \end{itemize}
  \item Simple
    \begin{itemize}
    \item Use HW-registers
    \item Perform \alert{peephole} optimizations\\
      (substituting instruction sequences with more efficient ones)
    \end{itemize}
  \item Advanced
    \begin{itemize}
    \item Run various analyses on code
    \item Generate multiple instances of same method for different cases
    \item \ldots
    \end{itemize}
  \end{itemize}
\end{frame}

\begin{frame}{Interesting Facts About JIT}
  \begin{itemize}  \setlength{\itemsep}{\fill}
  \item Term originates from just-in-time manufacturing\\
    (aka kanban method)
  \item Appeared around the time of Java’s uptake
  \item Terrible misnomer (should be \alert{just too late})
  \item Universally misapplied\\
    (e.g., to dynamic compilation \alert{after} first execution)
  \item ``JIT'' is \alert{not a noun}
  \end{itemize}
\end{frame}

\section{When to JIT}

\begin{frame}{So! When should I JIT compile?}
  \begin{columns}
    \begin{column}{.4\linewidth}
      Key factors:
      \begin{itemize}  \setlength{\itemsep}{\fill}
      \item Speed of interpretation
      \item Speed of compilation
      \item Speed of compiled code
      \end{itemize}
    \end{column}
    \begin{column}{.5\linewidth}
      \centering
      \begin{tikzpicture}
        \begin{axis}[
          xlabel={Number of Executions},
          ylabel={Time},
          ylabel near ticks,
          ymin=0, xmin=0,
          ytick=\empty,
          axis y line=left,
          axis x line=bottom,
          small]
          % use TeX as calculator:
          \addplot[
          color=myyellow,
          mark=*, thick] coordinates { (-1,-1) (1,1) (2,2) (3,3) (4,4) (5,5) };
          \addplot[
          color=YellowGreen,
          mark=*, thick] coordinates { (-1,-2.5) (1,2.5) (2,3) (3,3.5) (4,4) (5,4.5) };
        \end{axis}
      \end{tikzpicture}
    \end{column}
  \end{columns}
\end{frame}

\begin{frame}{Measuring Speed}
  \begin{enumerate}  \setlength{\itemsep}{\fill}
  \item Measure the \alert{whole run} in \alert{clock cycles}
  \item \alert{Subtract} GC time, native code, and anything unrelated
  \item Divide by number of bytecodes to obtain \alert{cycles per bytecode}
  \end{enumerate}
\end{frame}

\begin{frame}{Measuring Speed}
  \begin{itemize}  \setlength{\itemsep}{\fill}
  \item Interpretation time \alert{$t_i$}
  \item Compilation time \alert{$t_c$}
  \item Execution time of compiled code \alert{$t_e$}
  \item Compiler over Interpreter ratio \alert{$\tau = {t_c \over t_i}$} (mnemonic: translate)
  \item Interpreter over Compiled code ratio \alert{$\rho = {t_i \over t_e}$} (mnemonic: run)
  \end{itemize}
\end{frame}

\begin{frame}{The break-even point}
  \centering
  \huge
  \alert{$\beta = {t_c \over t_i - t_e} = {{\rho \tau} \over \rho - 1} $}

  \vfill
  \normalsize
  Example: $t_e = 1$, $t_i = 2$, $t_c = 10$, $\tau = 5$, $\rho = 2$, $\beta = 10$
\end{frame}

\begin{frame}{log $\beta = {{\rho \tau} \over \rho - 1}$ for $t = 1, 2, 5, 10, 20, 50, 100$}
  \centering
  \begin{tikzpicture}
    \begin{axis}[
      width=.9\columnwidth,
      height=\axisdefaultheight,
      domain=0.001:20,
      xlabel={$\rho$},
      ylabel={$\beta$},
      ymode=log,
      ymin=0,
      xmin=0,
      axis y line=left,
      axis x line=bottom,
      samples=500]
      % use TeX as calculator:
      \addplot[color=myyellow!10, thick] { (100*x) / (x-1) }
      node[above,xshift=-1cm, yshift=-1mm] {\footnotesize$\tau = 100$};
      \addplot[color=myyellow!25, thick] { (50*x) / (x-1) }
      node[above,xshift=-1cm, yshift=-1mm] {\footnotesize$\tau = 50$};
      \addplot[color=myyellow!40, thick] { (20*x) / (x-1) }
      node[above,xshift=-1cm, yshift=-1mm] {\footnotesize$\tau = 20$};
      \addplot[color=myyellow!55, thick] { (10*x) / (x-1) }
      node[above,xshift=-1cm, yshift=-1mm] {\footnotesize$\tau = 10$};
      \addplot[color=myyellow!70, thick] { (5*x) / (x-1) }
      node[above,xshift=-1cm, yshift=-1mm] {\footnotesize$\tau = 5$};
      \addplot[color=myyellow!85, thick] { (2*x) / (x-1) }
      node[above,xshift=-1cm, yshift=-1mm] {\footnotesize$\tau = 2$};
      \addplot[color=myyellow, thick] { (x) / (x-1) }
      node[above,xshift=-1cm, yshift=-1mm] {\footnotesize$\tau = 1$};
    \end{axis}
  \end{tikzpicture}
\end{frame}

\begin{frame}{log $\mathrm{Speedup} = {{n t_i} \over t_c + n t_e} = {{n \rho} \over n + \rho \tau}$ for $\rho = 5$, $\tau = 100$}
  \centering
  \begin{tikzpicture}
    \begin{axis}[
      width=.9\columnwidth,
      height=6cm,
      domain=0:1000,
      xlabel={Number of Executions},
      ylabel={Speedup},
      ymode=log,
      axis y line=left,
      axis x line=bottom,
      samples=500]
      % use TeX as calculator:
      \addplot[color=myyellow, thick, dashed, samples=2] {1}
      node[below, xshift=-1cm] {$y = 1$};
      \addplot[color=myyellow, thick, dashed, mark=none]
      coordinates { (125,0.02) (125, 10) } node[right, yshift=-2mm] {$x = 125$};
      \addplot[thick] { (5*x) / (x + 5*100) };
    \end{axis}
  \end{tikzpicture}
\end{frame}

\section{How to JIT}

\begin{frame}{Interaction of JITed and non-JITed Code}
  \begin{itemize}  \setlength{\itemsep}{\fill}
  \item JITed code can invoke non-JITed code
  \item JITed code can fall-back to non-JITed counterpart
  \end{itemize}
\end{frame}

\begin{frame}{JITed Code Management}
  \begin{itemize}  \setlength{\itemsep}{\fill}
  \item Maintain a \alert{code cache} holding the JITed code
  \item VMs manage it to ensure new code can be added when needed
  \item Compilers rely on a big buffer to produce the code
  \item The JITed code is then copied to the code cache
  \item Calls to that code are redirected (e.g. through method tables)
  \end{itemize}
\end{frame}

\begin{frame}{When to Remove Code}
  \begin{itemize}  \setlength{\itemsep}{\fill}
  \item When the corresponding code \alert{becomes unreachable}\\
    (e.g. class unloading)
  \item When the code cache doesn't have enough \alert{space for new code}
  \end{itemize}
\end{frame}

\begin{frame}{What if Still Active?}
  \begin{itemize}  \setlength{\itemsep}{\fill}
  \item Keep an \alert{activation-counter}
  \item \alert{Scan} the stacks for activation records
  \end{itemize}
  \pause
  \vfill
  \centering
  Sounds familiar?
\end{frame}

\begin{frame}{Stack Scanning}
  \begin{itemize}  \setlength{\itemsep}{\fill}
  \item \alert{Walk the stack} looking for return addresses in the code
  \item The \alert{calling convention} must allow for stack scanning
    \begin{itemize}
    \item Standard placement of return addresses in the stack
    \item Access to saved registers (e.g. SP, FP) of suspended threads
    \end{itemize}
  \item Depends on underlying, architecture, OS, ad calling convention
  \item The OS can make this \alert{impossible}\\
    (e.g., by saving register state in kernel space).
  \end{itemize}
\end{frame}

\begin{frame}{What if Still Active?}
  \begin{itemize}  \setlength{\itemsep}{\fill}
  \item Let it be
  \item Drop it and patch any activation records pointing to it
    \begin{itemize}
    \item Do the housekeeping (fix state, fix return address, etc.)
    \item Make it resume in the interpreter
    \item Called \alert{dynamic deoptimization}
    \item It's slow!!!
    \end{itemize}
  \end{itemize}
\end{frame}

\begin{frame}{What About Fragmentation in the Code Cache}
  \begin{itemize}  \setlength{\itemsep}{\fill}
  \item Self VM uses \alert{compaction}
    \begin{itemize}
    \item Similar to GC
    \item Stop the world
    \item Compact code cache
    \item Update links (method tables) and return addresses (activation records)
    \end{itemize}
  \item HotSpot\texttrademark{} drops code till enough space becomes available
  \end{itemize}
\end{frame}

\begin{frame}{Using the Heap}
  Make compiled code an object in the heap

  \vfill

  Pros \tik
  \begin{itemize}
  \item \alert{Share} memory management mechanisms
  \end{itemize}

  \vfill

  Cons \ex
  \begin{itemize}
  \item Compiled code usually \alert{lives long}
  \item A tiny bug may give \alert{write access to generated code}
  \end{itemize}
\end{frame}

\section{Impact of JITed Code on Garbage Collection}

\begin{frame}{Impact of JITed Code on Garbage Collection}
  \begin{itemize}  \setlength{\itemsep}{\fill}
  \item \alert{More} roots
  \item Pointers might be:
    \begin{itemize}
    \item in \alert{HW-registers}
    \item \alert{stack frames} of JITed code
    \item or even \alert{hardcoded} in the JITed code (as literals)
    \end{itemize}
  \item A HW-register might even contain a pointer in the middle of (de)construction, e.g. pointer arithmetic
  \end{itemize}
\end{frame}

\begin{frame}{Keeping Track of Roots in Registers}
  \begin{itemize}  \setlength{\itemsep}{\fill}
  \item All activations (except the top) are at a \alert{call site}
  \item For each call site keep a \alert{register map}
  \item The register map indicates which registers are live
  \end{itemize}
\end{frame}

\begin{frame}{Keeping Track of Roots in Registers: Example}
  \begin{itemize}  \setlength{\itemsep}{\fill}
  \item Self compilers inject a register map word after each call:\\
    each bit represents a register, with the bit set if the register is live and has an oop
  \item The return sequence skips over the map
  \item A simple stack scan locates all the register maps
  \end{itemize}
\end{frame}

\begin{frame}{Keeping Track of Roots in Registers of Top Frame}
  \begin{itemize}  \setlength{\itemsep}{\fill}
  \item No root call site with register map
  \item Can be suspended at any instruction
  \end{itemize}

  \vfill
  \hrule
  \vfill

  \begin{itemize}  \setlength{\itemsep}{\fill}
  \item Only allow suspension at GC \alert{safe points}, with register maps
  \item Maintain register maps at fixed locations and use \alert{abstract interpretation} to derive the map at current point
  \item \alert{Replay} the compiler to produce the register map(s)
  \end{itemize}
\end{frame}

\begin{frame}{GC Safe Points}
  Safe points are placed at JITed \alert{code entry and back-branches}
  \vfill
  At safe points \alert{ask the VM} whether the thread should suspend
  \vfill
  To stop-the-world all threads need to \alert{reach a safe point}
  \vfill
  Waiting others to reach a safe point, threads can \alert{scan themselves}
\end{frame}

\begin{frame}{Keeping Track of Roots on the Stack}
  Keep \alert{stack maps} for each call site
  \vfill
  Keep the stack maps in the code cache
  \vfill
  Add a pointer to them after the call
\end{frame}

\begin{frame}{Keeping Track of Roots in the Code}
  If the compiler hardcodes pointers in the JITed code:
  \vfill
  \begin{itemize}  \setlength{\itemsep}{\fill}
  \item It might be in complex form, e.g. \alert{piecemeal assembly}
  \item The compiler emits a table identifying the locations of these refs
  \item The VM needs to be able to use them as roots
  \item The VM needs to be able to alter them if needed
  \end{itemize}
\end{frame}

% Closing slide
% \subsection{Closure}
% % {
% % \setbeamertemplate{footline}{}
% % \begin{frame}[noframenumbering]
% %  \frametitle{\fontspec{Purisa}\textbf{Thank You!}}
% %  \centering
% %  \titlepage
% % \end{frame}
% % }
% \begin{frame}
%   % \centering
%   \Huge
%   \alert{Can managed programming languages run on future
%   many-core architectures?}
%   % \vspace{-2em}
%   \vfill
%   \begin{flushright}
%     \fontspec{Purisa}\structure{Thank You!}
%   \end{flushright}
% \end{frame}

%   BACKUP SLIDES

\end{document}
